\section{方法}

\subsection{语汇核验}

对候选词 $q$，定义完全匹配次数
\begin{equation}
H(q)=\sum_{j=1}^{|T|}\mathbf{1}[T_j=q],
\end{equation}
其中 $T$ 为原书字符序列。该量只回答原句是否存在及出现多少次。本文引用的“这件工具的名字叫做‘知行合一’”“此学即所谓‘致良知’之心学，传我此学者，名王守仁”均先经同一接口逐字核验；前一关键词命中12次，后一关键词命中2次。

\subsection{关系类别构成}

给定人物集合 $G$，取至少一个端点属于 $G$ 的关系集合 $R_G$。类别 $c$ 的构成比为
\begin{equation}
P(c\mid G)=\frac{|\{r\in R_G:r.category=c\}|}{|R_G|}.
\end{equation}
思想相关组由结构化记录明确连接到心学、王学、讲学或书院的10人组成；权力参照组取张居正、严嵩。传播替代量为“师生同门”边占比；权力接合替代量为君臣从属、政治攻讦、敌对冲突与同盟友党四类之和。两组人数不同，因此只比较占比，不比较边数绝对值。

主分析删除章节标记为“推导”的同事件共现边。稳健性分析恢复全部边：若组间方向保持，说明结论不依赖是否保留自动共现。另报告王守仁、徐阶、顾宪成、高攀龙、张居正、严嵩六个锚点的度数和显式类别，防止组平均掩盖角色差异。

\subsection{时间序列}

思想事件按年份计数，并分为1480--1507、1508--1528、1529--1578、1579--1600四个窗口。窗口不是历史分期，而是把形成、传播间隔与晚期制度冲突分开显示的描述工具。任何“先传播、后冲突”的表述都只指该书的事件编码顺序，不构成因果推断。

全部统计由\texttt{\_stats.py}调用共用\texttt{mingdata}生成，输出\texttt{.dump/\_stats\_ideology.json}。两张图只读取该底稿，不在作图脚本中手填数字。
